% Textual recovery, 18 November 2026.
% Through P.16.12: restored checkpoint08 source; result wrapper promoted.
% P.16.13--16: fresh transcription of the separately reviewed post-checkpoint corollary.
% This is not claimed byte-identical to the lost final source.
\subsection{The exact Brauer--Manin criterion for the remaining weight-three match}

\begin{result}{P.16}{Proved}
{The exact central Brauer criterion and the scope of base-tripod Brauer classes}
Fix an actual integral pro-five class-two predecessor and a global
class-three lift $t$. At a set $\Sigma$ of finite primes different
from five, suppose the predecessor matches specified actual local
point sections. Retain the unique lower aligning frames of B.26.
Put
\[
 T=\mathbb Z_5(3)^b,\quad M_n=T/5^nT,\quad
 M_n^D=\operatorname{Hom}(M_n,\mu_{5^n})
       =(\mathbb Z/5^n)(-2)^b,\quad S_5=\{5,\infty\}.
 \tag{P.16.1}
\]
For the tripod $b=b_3=2$. Write $d_p\in H^1(\mathbb Q_p,T)$
for the required correction to $t$ and $d_{p,n}$ for its reduction.
The existence of one integral global correction matching every
$p\in\Sigma$ is equivalent to the following condition:
for every $n$, the classes $d_{p,n}$ admit a completion
$(a_{v,n})_v$ such that
\[
 \begin{gathered}
 a_{p,n}=d_{p,n}\quad(p\in\Sigma),\qquad
 a_{v,n}\in H^1_{\rm nr}(\mathbb Q_v,M_n)
                         \quad(v\ne5,\infty),\\
 \sum_v\operatorname{inv}_v(a_{v,n}\cup c_v)=0
                \quad\hbox{for every }c\in H^1(\mathbb Q,M_n^D).
 \end{gathered}
 \tag{P.16.2}
\]
At five the completion is unrestricted; at the real place its
degree-one class is zero. At places outside $\Sigma$ these are
choices of correction classes, not assertions that a local point
has the prescribed predecessor. The sum in P.16.2 is finite for
each $c$. All classes $d_{p,n}$ are unramified, even when
$\Sigma$ is infinite.

P\'al--Schlank's theorem applies to this completion problem with
the actual predecessor retained. Its extra Brauer condition is
exactly the central matching obstruction of E.212. In particular,
every one-place or two-place failure there gives a nonbase
algebraic Brauer class detecting that failure, independently of
all choices at five and at the unspecified places. This comparison
does not prove that P.16.2 holds for the ordinary pair produced
in E.213.

\end{result}

\begin{proof}
\emph{The primary theorem is an exact image assertion.}
For a finite embedding problem $\mathcal E$ over a global field
$F$, P\'al--Schlank define $\operatorname{Sol}_{\mathbb A}(\mathcal E)$
using local solutions unramified at almost every place.
Solutions are taken up to conjugacy by the finite kernel; they
need not be surjective. If the kernel is abelian and its order
is prime to $\operatorname{char}(F)$, their Theorem 1.2 states
\[
 r\bigl(\operatorname{Sol}(\mathcal E)\bigr)
  =\operatorname{Sol}_{\mathbb A}^{\operatorname{Br}}(\mathcal E)
  =\operatorname{Sol}_{\mathbb A}^{\operatorname{Br}_1}(\mathcal E).
 \tag{P.16.3}
\]
There is no closure on the left. We use the opened author
version, whose statement is on p.~3 and whose proof is in
\S8, pp.~27--28
\citep[Theorem 1.2; \S1, pp.~1--3; \S8, pp.~27--28]{LitPalSchlank2022}.
Its cohomological form, Theorem 5.4 on pp.~14--15, says that
the image of global first cohomology is the annihilator of the
global dual first cohomology in the restricted adelic product.
Proposition 8.1 identifies this pairing with the algebraic
Brauer pairing. These statements retain the chosen embedding
problem, rather than replacing its predecessor.

\emph{Apply the theorem to the correction problem.}
Let $C_n$ be the image of $\chi_5^3$ in
$(\mathbb Z/5^n)^\times$. Consider the finite split problem
\[
 G_{\mathbb Q}\xrightarrow{\chi_5^3\bmod5^n}C_n,\qquad
 M_n\rtimes C_n\longrightarrow C_n .
 \tag{P.16.4}
\]
Its solution set is $H^1(\mathbb Q,M_n)$, and its local
solution sets are the corresponding local first-cohomology
groups. This auxiliary split problem parametrizes changes
to the fixed lift $t$. It does not replace the actual
nonabelian predecessor by its cyclotomic quotient: an integral
solution will be added centrally to $t$ itself.

The integral tame relation gives
$H^1(\mathbb Q_p,T)=H^1_{\rm nr}(\mathbb Q_p,T)$ for
$p\ne5$, so every $d_{p,n}$ is unramified. Consequently any
completion in P.16.2 is an adelic solution of P.16.4. No finite
support condition on its Frobenius values is required by
the adelic definition. Only ramification must be absent
almost everywhere; here it is absent everywhere outside five.
The first line of P.16.2 permits all unramified classes at an
unspecified prime. Requiring the zero correction there would
be an additional condition and could change the answer.

By P.16.3, a completion satisfying the second line of P.16.2
is the localization of some
$h_n\in H^1(\mathbb Q,M_n)$. Since every local class outside
five is unramified, $h_n$ factors through $G_{\mathbb Q,S_5}$.
Indeed the kernel of $G_{\mathbb Q}\to G_{\mathbb Q,S_5}$
is normally generated by those inertia groups; the coefficient
action is already unramified there, and the restricted
cocycle vanishes on them. Conversely, a class in
$H^1(G_{\mathbb Q,S_5},M_n)$ supplies a completion satisfying
P.16.2 by global reciprocity.

\emph{Recover one integral correction without changing the source.}
E.209 and the coefficient exact sequence identify
\[
 H^1(G_{\mathbb Q,S_5},M_n)
   =H^1(G_{\mathbb Q,S_5},T)/5^n
   =(\mathbb Z/5^n)^b .
 \tag{P.16.5}
\]
For each $n$, let $K_n\subseteq\mathbb Z_5^b$ be the set
of integral correction parameters whose local reductions
match every $d_{p,n}$. Equation P.16.5 shows that P.16.2 is
equivalent to $K_n\ne\varnothing$. Each $K_n$ is closed:
it is the intersection of the inverse images of prescribed
classes under continuous maps to finite groups.
Also $K_{n+1}\subseteq K_n$. Compactness of $\mathbb Z_5^b$
therefore gives an element of $\bigcap_nK_n$.
Its twist of $t$ is the desired actual integral lift.
The converse follows by reducing one successful integral
correction. Thus this passage requires no choice of compatible
solutions on varying arithmetic fields or with new ramification.
It is a conclusion at this one fixed central lifting step.

\emph{Which Brauer classes remain after the base classes vanish.}
For P.16.4 the algebraic Brauer group is naturally
$H^1(\mathbb Q,M_n^D)$, through the injective edge map of
\citep[Lemma 2.2 and Definition 2.3, author pp.~5--6]{LitPalSchlank2022}.
This is a group modulo the arithmetic base Brauer group.
An explicit representative uses the fibre product
$M_n\rtimes G_{\mathbb Q}$ and its universal central cocycle
$u(m,g)=m$. For a global dual cocycle $c$, the cup product
$u\cup c$, followed by $\mu_{5^n}\hookrightarrow\overline{\mathbb Q}^{\,*}$,
is zero on the reference section. Its edge class is $c$
up to the uniform sign convention for that edge map.
These classes run through the algebraic Brauer group, and
their evaluation on a correction $a$ is
\[
          \sum_v\operatorname{inv}_v(a_v\cup c_v).
 \tag{P.16.6}
\]
This is also the normalization supplied by Proposition 8.1,
author pp.~27--28. In particular a nonzero edge class is
not the pullback of an arithmetic base class. The already
global abelian pair gives no automatic vanishing of P.16.6
in this argument. To infer such vanishing from geometric
Brauer groups would require an additional comparison of
these classes and of the actual lifted local points.
No such comparison is assumed here.

\emph{The exact dual class of a two-place mismatch.}
It suffices first to consider one Hall coordinate. Set
$A_n=(\mathbb Z/5^n)(3)$ and $D_n=A_n^D=(\mathbb Z/5^n)(-2)$.
The negative twist must be retained; $-2$ can be replaced by
$2$ only modulo five, not modulo arbitrary $5^n$.
For $S_5=\{5,\infty\}$ we have
\[
 H^1(G_{\mathbb Q,S_5},D_n)
       =H^2(G_{\mathbb Q,S_5},D_n)=0 .
 \tag{P.16.7}
\]
For $n=1$, E.209 gives the second equality for the actual
negative twist. Global Euler characteristic gives the first:
the invariants vanish, complex conjugation acts trivially,
and the second cohomology vanishes. The successive coefficient
sequences prove the first equality for every $n$; the second
also follows directly from E.209. This uses precisely the
finite global Euler formula used in E.210
\citep[I.5.1, p.~67]{CohMilneADT}.

Let $F\subseteq\Sigma$ be finite. Finite purity and P.16.7
give an isomorphism
\[
 H^1(G_{\mathbb Q,S_5\cup F},D_n)
  \xrightarrow{\ \sim\ }
  \bigoplus_{p\in F}H^0(\mathbb F_p,D_n(-1)).
 \tag{P.16.8}
\]
Local Tate duality identifies the summand at $p$ with the
dual of $H^1_{\rm nr}(\mathbb Q_p,A_n)$. One can define the
residue identification by this pairing, fixing its sign.
Thus any additive character
\[
 L:\prod_{p\in F}H^1_{\rm nr}(\mathbb Q_p,A_n)
                         \longrightarrow (5^{-n}\mathbb Z)/\mathbb Z
 \tag{P.16.9}
\]
determines a unique class $c_L$ by its residues in P.16.8.

Choose the integral generator $\zeta_3$ of E.212, with
local values $\tau_p$. Suppose $L$ annihilates the localization
of the reductions of $\mathbb Z_5\zeta_3$. Reciprocity gives
\[
 0=\sum_v\operatorname{inv}_v
          \bigl((\zeta_3\bmod5^n)_v\cup(c_L)_v\bigr)
   =\operatorname{inv}_5
          \bigl((\zeta_3\bmod5^n)_5\cup(c_L)_5\bigr).
 \tag{P.16.10}
\]
Outside $F\cup\{5\}$ both classes are unramified; the real
term is zero. At five, $H^0$ and $H^2$ for $A_n$ and $D_n$
vanish. More precisely, local duality gives
$H^2_{\rm cts}(\mathbb Q_5,\mathbb Z_5(3))=0$, because
$(\mathbb Q_5/\mathbb Z_5)(-2)$ has no invariants.
The coefficient sequence and E.210's integral rank-one
calculation therefore give
\[
 H^1(\mathbb Q_5,A_n)
   =H^1_{\rm cts}(\mathbb Q_5,\mathbb Z_5(3))/5^n
   \simeq\mathbb Z/5^n.
\]
By E.210, $(\zeta_3\bmod5^n)_5$ generates this cyclic group.
Perfect local Tate duality then makes $H^1(\mathbb Q_5,D_n)$
its cyclic dual. Thus cyclicity is obtained from the integral
coefficient sequence, not from a cardinality calculation alone.
Equation P.16.10 therefore implies
\[
                         (c_L)_5=0.
 \tag{P.16.11}
\]
For every completion of the prescribed support discrepancies
by unramified classes away from five, P.16.6 evaluated at $c_L$
is consequently exactly $L((d_{p,n})_{p\in F})$.
All unspecified places and the coefficient place contribute
zero. This remains true coordinatewise for $b$ Hall coordinates.

If a finite tuple is outside the localization image, finite
Pontryagin duality supplies a character $L$ annihilating that
image and nonzero on the tuple. E.212 permits $|F|\le2$ after
choosing a sufficiently large coefficient exponent whenever
integral matching fails. Hence $c_L$ is an explicit kind of
nonbase algebraic Brauer witness for precisely that failure.

For example, suppose $p^3\equiv q^3\equiv1\pmod{5^n}$ and
use the common local frames. Taking the residues
$\tau_q$ at $p$ and $-\tau_p$ at $q$ gives the evaluation
\[
       \frac{\tau_q d_{p,n}-\tau_p d_{q,n}}{5^n}
                         \quad\hbox{in }\mathbb Q/\mathbb Z.
 \tag{P.16.12}
\]
When both periods are units, this detects their incompatible
normalized ratios. In nonunit cases one uses the full
annihilator in P.16.9; the single displayed cross product
need not detect an individually impossible divisibility
condition. This is exactly the distinction retained in
E.212.5--6.

The theorem therefore applies, but does not make the remaining
pairings vanish. Ordinary abelian local matching and the
existence of the actual class-three lift do not constitute
a proof of that extra vanishing. The exact unresolved
condition remains E.212 for the actual lift and actual
point paths, followed by the analogous conditions at
subsequent weights.


\paragraph{Base-tripod Brauer classes do not supply the extra central condition.}
Put $U=\mathbb P^1_{\mathbb Q}\setminus\{0,1,\infty\}$ and use
$x$ for its standard coordinate. There is a split isomorphism
\[
 \operatorname{Br}(U)
 =\operatorname{Br}(\mathbb Q)
   \oplus H^1(\mathbb Q,\mathbb Q/\mathbb Z)
   \oplus H^1(\mathbb Q,\mathbb Q/\mathbb Z),
 \qquad (C,\alpha,\beta)\longmapsto
 C+(\alpha,x)+(\beta,1-x).
 \tag{P.16.13}
\]
Here $(\alpha,f)$ denotes the cyclic Brauer class defined by the
character $\alpha$ and the unit $f$. Consequently
\[
 \operatorname{Br}(U)[5^n]
 =\operatorname{Br}(\mathbb Q)[5^n]
   \oplus H^1(\mathbb Q,\mathbb Z/5^n)
   \oplus H^1(\mathbb Q,\mathbb Z/5^n).
 \tag{P.16.14}
\]
These are assertions about the base tripod. They do not describe
the Brauer groups of its finite covers or identify the nonbase
embedding-problem classes in P.16.6 with classes on $U$.

To prove P.16.13, inject the Brauer group of the regular curve into
its function-field Brauer group. The residue sequence identifies
its image with classes unramified at every closed point of $U$.
The Faddeev exact sequence for $\mathbb Q(x)$ then leaves only
residues at $0,1,\infty$, whose sum is zero. The two displayed
symbols have residues $(\alpha,0,-\alpha)$ and
$(0,\beta,-\beta)$, respectively. Subtracting them leaves a
constant class; taking their residues also proves directness.
We use \citet[Theorem~6.9.1, p.~176]{RecoveryGilleSzamuely2006}
and \citet[Theorem~1.5.2, p.~44; Theorem~3.5.4, pp.~88--89;
Proposition~3.6.1, pp.~89--91]{RecoveryColliotTheleneSkorobogatov2019}.
Theorem~1.5.2 is the Faddeev theorem; it is not being cited as a
purity theorem. The geometric Brauer vanishing underlying the
curve calculation is also supplied by Tsen's theorem
\citep[Theorem~6.2.8, p.~143]{RecoveryGilleSzamuely2006}.
Taking $5^n$-torsion in this split decomposition proves P.16.14.

Now let $(A,B)\in\widehat{\mathbb Q^{\times}}^{\,2}$ be an
ordinary completed pair. For every finite prime $p$, choose
$x_p\in U(\mathbb Q_p)$ whose completed Kummer coordinates are
the localizations of $(A,B)$. The tuple need not be an adelic
point of $U$. Fix $n$ and a class
$\mathcal B=C+(\alpha,x)+(\beta,1-x)\in\operatorname{Br}(U)[5^n]$.
Choose rational representatives $a_n,b_n\in\mathbb Q^{\times}$
of $A,B$ modulo $5^n$-th powers. Such representatives exist by
the definition of the completion. Since a cyclic symbol of
order dividing $5^n$ depends only on that power class,
\[
 \mathcal B(x_p)
   =\bigl(C+(\alpha,a_n)+(\beta,b_n)\bigr)|_{\mathbb Q_p}
   \quad\text{for every finite }p.
 \tag{P.16.15}
\]
The expression in parentheses is one global Brauer class.
Its local invariants therefore have finite support and sum to
zero. Its real invariant is zero because its order divides the
odd integer $5^n$. Thus
\[
 \#\{p:\operatorname{inv}_p\mathcal B(x_p)\ne0\}<\infty,
 \qquad
 \sum_{p<\infty}\operatorname{inv}_p\mathcal B(x_p)=0.
 \tag{P.16.16}
\]
The global reciprocity input is
\citet[Theorem~12.1.8, p.~243]{RecoveryColliotTheleneSkorobogatov2019}.
This argument establishes the finite sum directly and does not
assume adelic integrality of the original tuple. It supplies no
vanishing of P.16.6: the correction problem retains an actual
nonabelian predecessor and its point paths, and its nonbase
classes have not been compared with base-tripod Brauer classes.

\end{proof}
